\sectionheading{5.13 অ্যাসেম্বলার (Assembler)}

Assembly Language-এ Program লেখার সময় 0 ও 1-এর পরিবর্তে
\textenglish{\ttfamily ADD, SUB, MOV, JMP} ইত্যাদি
\textbf{Mnemonic Code} এবং প্রয়োজন অনুযায়ী Symbol বা Label
ব্যবহার করা হয়।

এসব Instruction Programmer-এর জন্য সহজ হলেও Processor সরাসরি
বুঝতে পারে না। তাই Assembly Language-এ লেখা Program-কে
Machine Language-এ রূপান্তর করার জন্য \textbf{Assembler} প্রয়োজন।

%================================================
% Definition
%================================================
\booksubheading{Assembler কাকে বলে?}

\begin{definitionbox}

\textbf{\color{darkTeal}সংজ্ঞা}

যে Translator Program \textbf{Assembly Language-এ লেখা Program-কে
Machine Language বা Machine Code-এ রূপান্তর করে}, তাকে
\textbf{Assembler} বলে।

\end{definitionbox}

\textbf{সহজভাবে মনে রাখবে—}

\begin{center}

\textbf{
Assembly Language
\(\longrightarrow\)
Assembler
\(\longrightarrow\)
Machine Language
}

\end{center}

%================================================
% Need
%================================================
\booksubheading{Assembler কেন প্রয়োজন?}

Assembly Language-এ Mnemonic Instruction ব্যবহার করা হলেও CPU
মূলত Binary Machine Instruction বুঝতে পারে।

Assembler এসব Mnemonic ও Symbolic Instruction-কে সংশ্লিষ্ট
Processor-এর বোধগম্য Machine Code-এ রূপান্তর করে।

\begin{examplebox}{সহজ উদাহরণ}

ধরা যাক Assembly Program-এ লেখা আছে—

\begin{center}
{\textenglish{\ttfamily ADD A, B}}
\end{center}

এখানে \textenglish{\ttfamily ADD} যোগ করার Operation প্রকাশ করে।

Assembler এই Instruction-কে সংশ্লিষ্ট Processor-এর উপযোগী
Binary Machine Instruction-এ রূপান্তর করবে।

\end{examplebox}

%================================================
% Assembling
%================================================
\booksubheading{Assembling কী?}

\begin{definitionbox}

\textbf{\color{darkTeal}Assembling}

Assembly Language-এর Program-কে Assembler-এর সাহায্যে
Machine Code-এ রূপান্তরের প্রক্রিয়াকে \textbf{Assembling} বলে।

\end{definitionbox}

\begin{center}

\textbf{
Assembler = অনুবাদক Program
\qquad
Assembling = অনুবাদের প্রক্রিয়া
}

\end{center}

%================================================
% Flow
%================================================
\booksubheading{Assembler-এর কাজের ধারা}

\begin{center}
\begin{adjustbox}{max width=0.98\linewidth}
\fcolorbox{darkBlue}{softBlue}{
    \parbox[c][9mm][c]{0.22\linewidth}{
        \centering\textbf{Assembly\\Source Program}
    }
}
\quad
\(\longrightarrow\)
\quad
\fcolorbox{mainTeal}{softGreen}{
    \parbox[c][9mm][c]{0.17\linewidth}{
        \centering\textbf{Assembler}
    }
}
\quad
\(\longrightarrow\)
\quad
\fcolorbox{darkBlue}{softBlue}{
    \parbox[c][9mm][c]{0.22\linewidth}{
        \centering\textbf{Machine /\\Object Code}
    }
}
\quad
\(\longrightarrow\)
\quad
\fcolorbox{mainTeal}{softGreen}{
    \parbox[c][9mm][c]{0.15\linewidth}{
        \centering\textbf{Execution}
    }
}
\end{adjustbox}
\end{center}

%================================================
% Main Functions
%================================================
\booksubheading{Assembler-এর প্রধান কাজ}

Assembler সাধারণত—

\begin{itemize}

\item Mnemonic Instruction-কে সংশ্লিষ্ট
\textbf{Machine Instruction}-এ রূপান্তর করে।

\item Opcode ও Operand-এর তথ্য Machine-এর উপযোগীভাবে অনুবাদ করে।

\item Program-এ ব্যবহৃত \textbf{Label বা Symbolic Address}
শনাক্ত ও প্রয়োজনীয় Address-এর সঙ্গে সম্পর্ক স্থাপন করে।

\item Assembly Language-এর Syntax ও Instruction-সংক্রান্ত
Error শনাক্ত করতে সাহায্য করে।

\item প্রয়োজনীয় \textbf{Object/Machine Code} তৈরি করে।

\end{itemize}

%================================================
% Label and Symbol
%================================================
\booksubheading{Label বা Symbolic Address কী?}

Assembly Program-এ কোনো Instruction বা Memory Location সহজে
শনাক্ত করার জন্য অর্থবোধক নাম ব্যবহার করা যেতে পারে।
এ ধরনের নামকে \textbf{Label বা Symbolic Name} বলা হয়।

যেমন—

\begin{center}

{\textenglish{\ttfamily START \qquad LOOP \qquad NEXT}}

\end{center}

Programmer-এর জন্য এসব নাম Numeric Address মনে রাখার তুলনায় সহজ।
Assembler এসব Symbol-এর সঙ্গে প্রয়োজনীয় Address-এর সম্পর্ক নির্ধারণ করে।

\begin{rememberbox}

\textbf{\color{deepOrange}মনে রাখবে}

Label Programmer-এর জন্য সহজ প্রতীকী নাম।

Assembler সেই Label বা Symbol-কে প্রয়োজনীয় Machine Address-এর
সঙ্গে সম্পর্কিত করে।

\end{rememberbox}

%================================================
% Features
%================================================
\booksubheading{Assembler-এর প্রধান বৈশিষ্ট্য}

\begin{tabularx}{\linewidth}{|>{\bfseries}p{0.27\linewidth}|Y|}
\hline

\rowcolor{softBlue}
\multicolumn{1}{|c|}{\textbf{বৈশিষ্ট্য}}
&
\multicolumn{1}{c|}{\textbf{ব্যাখ্যা}}
\\
\hline

Translator Program
&
এটি একটি অনুবাদক Program।
\\
\hline

Source Language
&
Assembly Language।
\\
\hline

Target Language
&
Machine Language বা Machine Code।
\\
\hline

Mnemonic Translation
&
Mnemonic-কে Machine Instruction-এ রূপান্তর করে।
\\
\hline

Symbol Handling
&
Label ও Symbolic Address নিয়ে কাজ করতে পারে।
\\
\hline

Error Detection
&
Assembly Syntax ও Instruction-এর ভুল শনাক্ত করতে সাহায্য করে।
\\
\hline

Machine Dependent
&
নির্দিষ্ট Processor-এর Instruction Set-এর সঙ্গে সম্পর্কিত।
\\
\hline

\end{tabularx}

%================================================
% Machine Dependent
%================================================
\booksubheading{Assembler কেন Machine Dependent?}

Assembly Language নির্দিষ্ট Processor Architecture-এর
\textbf{Instruction Set}-এর ওপর নির্ভরশীল।

এক Processor-এর Instruction, Register বা Mnemonic অন্য
Processor-এ একইভাবে নাও থাকতে পারে।

তাই Assembler-ও সাধারণত নির্দিষ্ট Processor Architecture-এর
Assembly Code অনুবাদের জন্য তৈরি হয়।

\begin{center}

\textbf{
Processor/Architecture পরিবর্তন
\(\longrightarrow\)
Instruction Set পরিবর্তন হতে পারে
\(\longrightarrow\)
Assembler-ও ভিন্ন হতে পারে
}

\end{center}

%================================================
% One Pass / Two Pass
%================================================
\booksubheading{One-pass ও Two-pass Assembler — সংক্ষিপ্ত ধারণা}

এটি মূল ধারণার তুলনায় কম গুরুত্বপূর্ণ; MCQ বা জ্ঞানমূলক
প্রশ্নের সুবিধার জন্য সংক্ষেপে জানলেই যথেষ্ট।

\begin{tabularx}{\linewidth}{|>{\bfseries}p{0.27\linewidth}|Y|}
\hline

\rowcolor{softBlue}
\multicolumn{1}{|c|}{\textbf{ধরন}}
&
\multicolumn{1}{c|}{\textbf{ধারণা}}
\\
\hline

One-pass Assembler
&
Source Program সাধারণত একবার Scan করে Translation সম্পন্ন
করার চেষ্টা করে।
\\
\hline

Two-pass Assembler
&
Source Program দুইবার Scan করে; প্রথম Pass-এ Symbol/Label
সম্পর্কিত তথ্য সংগ্রহ এবং পরবর্তী Pass-এ Machine Code তৈরির
কাজ করা হয়।
\\
\hline

\end{tabularx}

\begin{rememberbox}

\textbf{\color{deepOrange}Shortcut}

\textbf{One-pass = একবার Scan}

\textbf{Two-pass = দুইবার Scan}

\end{rememberbox}

%================================================
% Advantages
%================================================
\booksubheading{Assembler-এর সুবিধা}

\begin{compacttwocolpoints}

\item Programmer-কে সরাসরি 0 ও 1 দিয়ে Program লিখতে হয় না।

\item Mnemonic ব্যবহার করায় Program লেখা ও বোঝা সহজ হয়।

\item Label ও Symbolic Address ব্যবহার করা যায়।

\item Assembly Program-এর Error শনাক্ত করতে সহায়তা করে।

\item Hardware-এর কাছাকাছি Efficient Program তৈরিতে উপযোগী।

\end{compacttwocolpoints}

%================================================
% Limitations
%================================================
\booksubheading{Assembler-এর সীমাবদ্ধতা}

\begin{itemize}

\item নির্দিষ্ট Processor Architecture-এর ওপর নির্ভরশীল।

\item Assembly Programming-এর জন্য Processor ও Instruction Set
সম্পর্কে ধারণা প্রয়োজন।

\item High Level Language-এর তুলনায় বড় Program লেখা ও Maintain
করা কঠিন।

\item ভিন্ন Architecture-এ Program ব্যবহার করতে Code পরিবর্তনের
প্রয়োজন হতে পারে।

\end{itemize}

%================================================
% Example
%================================================
\booksubheading{Assembler-এর উদাহরণ}

Assembler Software-এর একটি পরিচিত উদাহরণ—

\begin{center}

\textbf{GAS — GNU Assembler}

\end{center}

%================================================
% Comparison
%================================================
\booksubheading{Compiler, Interpreter ও Assembler-এর পার্থক্য}

\begin{tabularx}{\linewidth}{
|>{\bfseries}p{0.17\linewidth}|Y|Y|Y|}
\hline

\rowcolor{softBlue}
\multicolumn{1}{|c|}{\textbf{বিষয়}}
&
\multicolumn{1}{c|}{\textbf{Compiler}}
&
\multicolumn{1}{c|}{\textbf{Interpreter}}
&
\multicolumn{1}{c|}{\textbf{Assembler}}
\\
\hline

Source Language
&
High Level
&
High Level
&
Assembly Language
\\
\hline

অনুবাদের ধরন
&
পুরো Program একসঙ্গে
&
সাধারণত Statement ধরে ধরে
&
Assembly Instruction অনুবাদ
\\
\hline

Object/Target Code
&
সাধারণত তৈরি করে
&
সাধারণত পৃথক Object Program তৈরি করে না
&
Machine/Object Code তৈরি করে
\\
\hline

Execution
&
Compilation-এর পরে
&
Translation-এর সঙ্গে ধাপে ধাপে
&
Machine Code তৈরির পরে
\\
\hline

প্রধান কাজ
&
High Level থেকে Target Code
&
High Level Code অনুবাদ ও Execute
&
Assembly থেকে Machine Code
\\
\hline

Error Handling
&
পুরো Program বিশ্লেষণ করে Error দেখাতে পারে
&
Error পাওয়া স্থানে সাধারণত থামে
&
Assembly Syntax/Symbol Error শনাক্ত করে
\\
\hline

\end{tabularx}

%================================================
% Shortcut
%================================================
\begin{rememberbox}

\textbf{\color{deepOrange}দ্রুত মনে রাখার কৌশল}

\begin{center}

\textbf{
Compiler
\(\rightarrow\)
High Level
\(\rightarrow\)
পুরো Program
}

\vspace{0.5mm}

\textbf{
Interpreter
\(\rightarrow\)
High Level
\(\rightarrow\)
Statement ধরে ধরে
}

\vspace{0.5mm}

\textbf{
Assembler
\(\rightarrow\)
Assembly
\(\rightarrow\)
Machine Code
}

\end{center}

\end{rememberbox}

%================================================
% Final Reminder
%================================================
\begin{rememberbox}

\textbf{\color{deepOrange}মনে রাখবে}

\begin{itemize}

\item Assembler একটি \textbf{Translator Program}।

\item এর Source Language হলো \textbf{Assembly Language}।

\item Target হলো \textbf{Machine Language/Machine Code}।

\item
\textbf{
Assembly Language
\(\rightarrow\)
Assembler
\(\rightarrow\)
Machine Code
}

\item Assembler-এর অনুবাদ প্রক্রিয়াকে
\textbf{Assembling} বলে।

\item Mnemonic, Label ও Symbolic Address নিয়ে Assembler কাজ করে।

\item Assembler \textbf{Machine Dependent}।

\item \textbf{GAS = GNU Assembler}।

\item One-pass ও Two-pass Assembler অতিরিক্ত ধারণা হিসেবে
সংক্ষেপে মনে রাখলেই যথেষ্ট।

\item Compiler ও Interpreter-এর সঙ্গে Assembler-এর সবচেয়ে
গুরুত্বপূর্ণ পার্থক্য হলো—
\textbf{Assembler-এর Source Language হচ্ছে Assembly Language}।

\end{itemize}

\end{rememberbox}

\small

%================================================
% 5.14 Interpreter
%================================================
